%% Question bank chapter file - generated by bank_tools.py build (do not edit while the bank is still being built:
%% it is regenerated from the per-test data in Question Bank\data). Every question is one bankq block:
%%   \begin{bankq}[id=..., source={...}, paper=.., qno=.., type=..., subtopic={...}, answer={...}, answer_src=key|solved,
%%                 working={...}, difficulty=1-10, diff_parts={K S I C F}, diff_note={...},
%%                 group=<passage id or empty>, also_in={duplicate copies in other tests}]
%% difficulty = min(10, K+S+I+C+F-1): K concepts 1-3, S steps 1-4, I insight 0-2, C computation 0-2, F format 0-1
%%   ...question LaTeX, options, figures (\bfig / \bfigi -> img/)...
%%   \end{bankq}
%% Passages shared by several questions are bankpassage blocks; their questions carry group=<passage id>.
\documentclass[10pt]{article}
\usepackage[a4paper,margin=1.8cm]{geometry}
\usepackage[T1]{fontenc}
\usepackage{amsmath,amssymb,mathtools}
\usepackage{graphicx}
\usepackage[export]{adjustbox}
\usepackage{xcolor,needspace,keyval}
\usepackage{enumitem}
\usepackage{booktabs,array,tabularx}
\usepackage[hidelinks]{hyperref}
\setlength{\parindent}{0pt}
\setlength{\parskip}{3pt}
\makeatletter
\define@key{bq}{id}{\def\bq@id{#1}}\define@key{bq}{source}{\def\bq@src{#1}}\define@key{bq}{paper}{\def\bq@paper{#1}}
\define@key{bq}{qno}{\def\bq@qno{#1}}\define@key{bq}{type}{\def\bq@type{#1}}\define@key{bq}{subtopic}{\def\bq@sub{#1}}
\define@key{bq}{answer}{\def\bq@ans{#1}}\define@key{bq}{answer_src}{\def\bq@asrc{#1}}\define@key{bq}{working}{\def\bq@work{#1}}
\define@key{bq}{group}{\def\bq@group{#1}}\define@key{bq}{also_in}{\def\bq@also{#1}}
\define@key{bq}{difficulty}{\def\bq@diff{#1}}\define@key{bq}{diff_parts}{\def\bq@dparts{#1}}\define@key{bq}{diff_note}{\def\bq@dnote{#1}}
\newenvironment{bankq}[1][]{\def\bq@work{}\def\bq@also{}\def\bq@group{}\def\bq@diff{?}\def\bq@dparts{}\def\bq@dnote{}\setkeys{bq}{#1}%
  \par\needspace{4\baselineskip}\medskip\noindent{\small\textbf{\bq@id}\hfill\color{gray}[\bq@src, Paper \bq@paper, Q\bq@qno]
  \bq@type\ $\cdot$ \bq@sub\ $\cdot$ \textbf{D\bq@diff}}\par\nopagebreak\relax}
  {\par{\small\color{teal}Answer: \textbf{\bq@ans}\ifx\bq@work\@empty\else\ (solved: \bq@work)\fi
  \ifx\bq@also\@empty\else\par\color{gray}Also in: \bq@also\fi}\par}
\define@key{bp}{id}{\def\bp@id{#1}}\define@key{bp}{source}{\def\bp@src{#1}}\define@key{bp}{questions}{\def\bp@qs{#1}}
\newenvironment{bankpassage}[1][]{\setkeys{bp}{#1}\par\needspace{4\baselineskip}\bigskip\noindent
  {\small\textbf{Passage \bp@id}\hfill\color{gray}[\bp@src; questions \bp@qs]}\par\nopagebreak\relax\itshape}{\par}
\makeatother
\newcommand{\bfig}[2]{\par\smallskip{\centering\includegraphics[width=#1pt,max width=0.62\linewidth,max height=0.32\textheight]{img/#2}\par}\smallskip}
\newcommand{\bfigi}[2]{\includegraphics[width=#1pt,max width=\linewidth,max height=0.32\textheight,valign=c]{img/#2}}
\newcommand{\bmissing}[1]{\textit{[figure missing in the source conversion: \detokenize{#1}]}}
\newenvironment{qoptions}{\begin{itemize}[leftmargin=1.8em,itemsep=1pt,topsep=2pt,label={}]}{\end{itemize}}
\newcommand{\opt}[2]{\item[(#1)] #2}
\providecommand{\fwstart}{}\providecommand{\fwend}{}\providecommand{\qfullsplit}{}
\newcolumntype{L}[1]{>{\raggedright\arraybackslash}p{#1}}
\newcolumntype{C}[1]{>{\centering\arraybackslash}p{#1}}
\newcolumntype{R}[1]{>{\raggedleft\arraybackslash}p{#1}}
\providecommand{\sectionhdr}[1]{}\providecommand{\subsectionhdr}[1]{}


\begin{document}
{\Large\bfseries Grignard --- Question Bank}\par\smallskip
{\small Chemistry chapter C16. Questions from all converted test papers, grouped by source series. Answers: printed key where it matches the paper, otherwise solved (answer\_src=solved, with working). Exact repeats across papers appear once, with also\_in listing the other copies.}\par\medskip\hrule

\section*{TOAS papers}
\begin{bankq}[id=TOAS04-P2-CHEM-S1-Q06, source={TOAS - Test 04 (15 Mar 2022)}, paper=2, qno=6, type=multi, subtopic={Grignard and organometallic reagents}, answer={C,D}, answer_src=key, difficulty=7, diff_parts={K2 S4 I1 C0 F1}, diff_note={Four Grignard or organometallic reactions: correct major product}]
Identify the reaction(s) with correct formation of major product ?
\begin{qoptions}
\opt{A}{\bfigi{153.9}{TOAS04_p2_chem_s1_q06_A.png}}
\opt{B}{\bfigi{78.2}{TOAS04_p2_chem_s1_q06_B.png}\bfigi{81.8}{TOAS04_p2_chem_s1_q06_3.png}}
\opt{C}{\bfigi{27.5}{TOAS04_p2_chem_s1_q06_4.png}$\xrightarrow[\mathrm{(2)\ H_2O}]{\mathrm{(1)\ CH_3{-}MgBr}}$\bfigi{63.4}{TOAS04_p2_chem_s1_q06_C.png}}
\opt{D}{\bfigi{143.3}{TOAS04_p2_chem_s1_q06_D.png}}
\end{qoptions}
\end{bankq}
\begin{bankpassage}[id=TOAS04-P2-CHEM-S3-G5-6, source={TOAS - Test 04 (15 Mar 2022)}, questions={5-6}]
\bfig{183.9}{TOAS04_p2_chem_s3_q04_1.png}
\end{bankpassage}
\begin{bankq}[id=TOAS04-P2-CHEM-S3-Q05, source={TOAS - Test 04 (15 Mar 2022)}, paper=2, qno=5, type=numerical, subtopic={Grignard and organometallic reagents}, answer={1}, answer_src=solved, working={Ph-MgBr + CO$_2$ gives benzoic acid B (DOU 5); LiAlH$_4$ gives benzyl alcohol C (DOU 4); $|n_1-n_2|=1$}, difficulty=4, diff_parts={K1 S3 I0 C0 F1}, diff_note={Grignard with CO2 then LiAlH4; difference of degrees of unsaturation}, group=TOAS04-P2-CHEM-S3-G5-6]
D.O.U. of compound B $=n_1$ and D.O.U. of compound C $=n_2$ then value of $|n_1-n_2|$
\end{bankq}
\begin{bankq}[id=TOAS04-P2-CHEM-S3-Q06, source={TOAS - Test 04 (15 Mar 2022)}, paper=2, qno=6, type=numerical, subtopic={Grignard and organometallic reagents}, answer={2}, answer_src=solved, working={(i) false (benzylic alcohol reacts immediately, not in 5 min); (ii) false (achiral); (iii) true: benzyl alcohol is an isomer of anisole; (iv) true: benzoic acid evolves CO$_2$; so 2}, difficulty=6, diff_parts={K2 S3 I1 C0 F1}, diff_note={Four statements about benzyl alcohol and benzoic acid}, group=TOAS04-P2-CHEM-S3-G5-6]
How many option(s) is/are correct\\
(i) (C) gives turbidity in 5 min with lucas reagent\\
(ii) (C) gives racemic mixture in product with $\mathrm{PCl_3}$\\
(iii) (C) is isomers of anisole\\
(iv) (B) gives $\mathrm{CO_2}$ gas with $\mathrm{NaHCO_3}$
\end{bankq}
\begin{bankq}[id=TOAS07-P1-CHEM-S4-Q07, source={TOAS - Test 07 (01 May 2022)}, paper=1, qno=7, type=numerical, subtopic={Grignard and organometallic reagents}, answer={7}, answer_src=key, difficulty=7, diff_parts={K2 S4 I1 C0 F1}, diff_note={Count esters, acid chlorides, carbonates needing 2 or more moles of RMgX}]
How many compounds out of following use two or more than two moles of Grignard reagent to convert into alcohol after complete hydrolysis.
\bfig{332.1}{TOAS07_p1_chem_s4_q07_1.png}
\end{bankq}
\begin{bankq}[id=TOAS07-P2-CHEM-S1-Q02, source={TOAS - Test 07 (01 May 2022)}, paper=2, qno=2, type=multi, subtopic={Grignard and organometallic reagents}, answer={A,B,C,D}, answer_src=key, difficulty=6, diff_parts={K2 S3 I1 C0 F1}, diff_note={Grignard addition then intramolecular substitution: ring size and ether}]
Which of the following is/are correct statement regarding major product?
\bfig{154.2}{TOAS07_p2_chem_s1_q02_1.png}
\begin{qoptions}
\opt{A}{Intramolecular nucleophilic substitution reaction takes place during the reaction sequence}
\opt{B}{5 membered ring is formed as major product}
\opt{C}{Grignard reagents ($\mathrm{CH_3MgBr}$) prefer to give nucleophilic addition on polar $\pi$-bond}
\opt{D}{Cyclic ether is formed as final major product}
\end{qoptions}
\end{bankq}
\begin{bankq}[id=TOAS09-P2-CHEM-S2-Q05, source={TOAS - Test 09 (30 May 2022)}, paper=2, qno=5, type=numerical, subtopic={Grignard and organometallic reagents}, answer={7}, answer_src=solved, working={1 mol for each acidic H (OH, SH), 2 for the lactone, 3 for the cyclic carbonate: 7}, difficulty=2, diff_parts={K1 S1 I0 C0 F1}, diff_note={Grignard and organometallic reagents: standard approach for this subtopic; 1 concept(s), 1 step(s)}]
Calculate total number of moles of Grignard (MeMgCl) required for complete reaction with following compound ?
\bfig{77.9}{TOAS09_p2_chem_s2_q05_1.png}
\end{bankq}
\end{document}
